\chapter{METHODOLOGY / SYSTEM DESIGN}
\label{ch:method}

This chapter explains \emph{what} was designed, \emph{why} it was designed in this way, and \emph{how} the experiments were controlled. The practical realisation of the design (code organisation, wire protocol, hardware rig, services and deployment) is described in Chapter~\ref{ch:implementation}.

\section{PROBLEM AND DESIGN OBJECTIVE}
\label{sec:problem}

The objective is a software layer above a BMS that converts voltage, current and time into health outputs an operator can act on, with two design constraints that shape everything below.

\begin{enumerate}
    \item \textbf{Fit nothing across cells.} Chapter~\ref{ch:results} shows that models fitted across cells do not retain skill under a change of test protocol. Both working estimators (SOH and RUL) are therefore expressed relative to the same cell's own early behaviour, so there is no training cohort to transfer from.
    \item \textbf{Refuse what cannot be justified.} Every output is validated, banded by its measured error, or refused with a stated cause. A refusal is a correct outcome that the system knows, so it is returned as a value, rendered, and reflected in the process exit code.
\end{enumerate}

\subsection{Research Variables}

\begin{itemize}
    \item \textbf{Independent variables:} the ohmic correction used (none, cycler resistance, resistance from the load step), the voltage window (fixed or learned), the reference size, the RUL threshold, and the dataset (CALCE, Oxford, NASA).
    \item \textbf{Dependent variables:} per-cell mean absolute SOH error and its median over cells, coverage (cells measured out of cells scored), rank correlation of resistance, RUL error indexed by true and by predicted distance to end of life, and error stratified by the consistency of readings.
    \item \textbf{Controlled variables:} the reference definition (a quantity a BMS can hold), the primary window and tolerances, fixed before each study ran, and the scoring rule. Choices made after seeing a result are labelled exploratory.
\end{itemize}

\section{SYSTEM REQUIREMENTS}
\label{sec:requirements}

The requirements below were fixed before the implementation and are the basis on which the results of Chapter~\ref{ch:results} are judged. Table~\ref{tab:requirements} lists them.

\begin{table}[htbp]
	\caption{Functional and Non-Functional Requirements}
	\centering
	\small
	\renewcommand{\arraystretch}{1.15}
	\setlength{\tabcolsep}{4pt}
	\begin{tabular}{|C{0.07\linewidth}|L{0.5\linewidth}|L{0.2\linewidth}|}
		\hline
		\textbf{ID} & \textbf{Requirement} & \textbf{Type} \\
		\hline
		R1 & Accept telemetry from a CAN bus, a serial rig and dataset files & Functional \\
		\hline
		R2 & Convert every source to one schema with documented units and sign convention & Functional \\
		\hline
		R3 & Estimate SOH, resistance and RUL from voltage, current and time & Functional \\
		\hline
		R4 & Refuse, with a stated reason, when the evidence is insufficient & Functional \\
		\hline
		R5 & Serve all results over an HTTP API and a web dashboard & Functional \\
		\hline
		R6 & Fit no model across cells for the shipped estimators & Design constraint \\
		\hline
		R7 & Show the provenance of every displayed value & Non-functional \\
		\hline
		R8 & Run end to end from a clean checkout with one command each for tests and results & Non-functional \\
		\hline
		R9 & Need no external service or physical hardware in continuous integration & Non-functional \\
		\hline
	\end{tabular}
	\label{tab:requirements}
\end{table}

\section{ASSUMPTIONS AND CONSTRAINTS}

The design rests on a small number of explicit assumptions. Telemetry is sampled at a known and reasonably uniform rate, so that charge can be obtained by integrating current over time. The current sign convention of each source is known and is corrected in the loader. A cell's own early discharges are representative of its undegraded state, which is why the reference is taken from them when the rated capacity is not declared. The service is a local bench and analysis tool: it binds to the loopback address by default, has no authentication and keeps its state in memory. These constraints are recorded because several of the results and several of the design choices follow from them.

\section{OVERALL SYSTEM FRAMEWORK}
\label{sec:framework}

BEACON ingests telemetry from three transports, converges them on one schema, and scores them through one shared function. The result, including every refusal, is served over an HTTP API to a React client (Fig.~\ref{fig:framework}).

\begin{figure}[htbp]
    \centering
    \resizebox{0.97\linewidth}{!}{%
    \begin{tikzpicture}[box/.style={draw, rounded corners=2pt, align=center, minimum height=9mm, text width=3.6cm, font=\small},
                        wide/.style={draw, rounded corners=2pt, align=center, minimum height=9mm, text width=11.4cm, font=\small},
                        arr/.style={-{Stealth[length=2.2mm]}, thick}]
        \node[box] (can) at (0,0) {CAN bus\\(vehicle, DBC file)};
        \node[box] (ser) at (4.5,0) {USB serial rig\\(ESP8266, INA219, LM35)};
        \node[box] (ds)  at (9,0) {Research datasets\\(CALCE, NASA, Oxford)};
        \node[wide] (frame) at (4.5,-2.3) {Unified telemetry frame: \texttt{cell\_id}, \texttt{test\_time\_s}, \texttt{voltage\_v}, \texttt{current\_a}, \texttt{temperature\_c}, \texttt{soc}};
        \node[wide, minimum height=22mm] (score) at (4.5,-5.0) {\texttt{score\_telemetry\_frame}\\[1mm]segment cycles $\rightarrow$ coulomb count $\rightarrow$ field SOH, resistance, RUL\\ $\rightarrow$ behaviour flags $\rightarrow$ risk $\rightarrow$ health index\\ $\rightarrow$ Guardian $\rightarrow$ digital twin};
        \node[box] (api) at (0,-7.9) {FastAPI\\(port 8000)};
        \node[box] (gw)  at (4.5,-7.9) {Express gateway\\(port 5000)};
        \node[box] (ui)  at (9,-7.9) {React client\\(port 5173)};
        \draw[arr] (can.south) -- (can.south |- frame.north);
        \draw[arr] (ser.south) -- (ser.south |- frame.north);
        \draw[arr] (ds.south)  -- (ds.south |- frame.north);
        \draw[arr] (frame.south) -- (frame.south |- score.north);
        \draw[arr] (api.north |- score.south) -- (api.north);
        \draw[arr] (api.east) -- (gw.west);
        \draw[arr] (gw.east) -- (ui.west);
    \end{tikzpicture}}
    \caption{Overall system framework: three transports, one unified frame, one scoring function, and the service tier.}
    \label{fig:framework}
\end{figure}

\subsection{Unified Scoring Function}

The function \texttt{score\_telemetry\_frame} contains everything downstream of acquisition: segmentation, coulomb counting, features, risk, health, RUL, the Guardian explanation and the digital twin. It was extracted from the CAN pipeline when serial ingestion was added, instead of being copied. A second segmentation or feature extractor would mean that a disagreement between a CAN run and a serial run over the same battery could not be traced to either the data or the code. With the tail shared, a difference in output is a difference in the telemetry.

\section{DATASETS}
\label{sec:data}

Four sources are used (Table~\ref{tab:datasets}). They were chosen because they differ in chemistry, format, temperature and protocol.

\begin{table}[htbp]
	\caption{Data Sources and Their Role}
	\centering
	\small
	\renewcommand{\arraystretch}{1.15}
	\setlength{\tabcolsep}{3pt}
	\begin{tabular}{|L{0.15\linewidth}|L{0.31\linewidth}|L{0.14\linewidth}|L{0.28\linewidth}|}
		\hline
		\textbf{Source} & \textbf{Cells and conditions} & \textbf{Role} & \textbf{Ground truth used} \\
		\hline
		CALCE CS2/CX2 & 22 prismatic LCO cells, room temperature, 7 protocols (constant-current, multi-rate, real partial cycling) & Development and primary validation & Cycler-measured discharge capacity \\
		\hline
		NASA PCoE & 18650 cells, 9 protocols, ambient 4--44~$^\circ$C; 31 cells in the benchmark, 27 in the cross-dataset study & Benchmark and third-dataset test & NASA's computed capacity \\
		\hline
		Oxford Dataset 1 & 8 Kokam NMC/LCO pouch cells, 740~mAh, 40~$^\circ$C, 519 characterisations & External validation, unchanged & 1C capacity at each characterisation \\
		\hline
		Bench rig & 1 HONGLI 18650 cell, at rest in the two real captures (83 and 85 records) & Data acquisition only & Frame integrity and timing \\
		\hline
	\end{tabular}
	\label{tab:datasets}
\end{table}

CALCE records span many files whose test time restarts per file. A monotonic time base is constructed by offsetting each restart, not by sorting timestamps, which would interleave overlapping segments. The CALCE set includes real partial cycling (CS2 Types 5 and 6): thousands of discharges of about 25\% depth in a fixed band, with a full capacity check roughly every 100 cycles. Type~6 cycles at the top of charge (about 4.07 to 3.78~V) and Type~5 near empty (about 3.69 to 2.70~V). The Oxford set additionally holds one real drive-cycle discharge of Cell~1 sampled at 1~Hz, used only for external validation.

\section{UNIFIED TELEMETRY SCHEMA}
\label{sec:schema}

Public datasets, vehicle frames and a serial rig describe the same measurements with different names, units and missing-value conventions, so every source is converted into one schema before any feature is computed. Table~\ref{tab:schema} lists the core columns.

\begin{table}[htbp]
	\caption{Core Columns of the Unified Schema}
	\centering
	\small
	\renewcommand{\arraystretch}{1.15}
	\setlength{\tabcolsep}{4pt}
	\begin{tabular}{|L{0.24\linewidth}|C{0.08\linewidth}|C{0.17\linewidth}|L{0.28\linewidth}|}
		\hline
		\textbf{Column} & \textbf{Unit} & \textbf{Required} & \textbf{Description} \\
		\hline
		\texttt{dataset} & text & Yes & Source dataset name \\
		\hline
		\texttt{cell\_id} & text & Yes & Cell, module or vehicle identifier \\
		\hline
		\texttt{cycle} or \texttt{timestamp} & -- & One of two & Ordering of the samples \\
		\hline
		\texttt{voltage\_v} & V & Yes & Terminal voltage \\
		\hline
		\texttt{current\_a} & A & Yes & Current with a documented sign \\
		\hline
		\texttt{temperature\_c} & $^\circ$C & Yes & Battery temperature \\
		\hline
		\texttt{capacity\_ah} & Ah & Recommended & Measured capacity \\
		\hline
		\texttt{soc} & \% & Recommended & State of charge, 0--100 \\
		\hline
		\texttt{resistance\_ohm} & $\Omega$ & Optional & Internal resistance if available \\
		\hline
	\end{tabular}
	\label{tab:schema}
\end{table}

\subsection{Derived Fields}

Three fields are derived from the core columns and regenerated after every preprocessing run. Power is the product of voltage and current. The operating mode is inferred from the current when the source does not give one: a current above $+0.02$~A is charge, below $-0.02$~A is discharge, and a current inside that band is rest. The state-of-charge band assigns each sample to one of four ranges, below 20\% (deep-discharge region), 20 to 80\% (preferred region), 80 to 90\% (elevated) and 90\% and above (high-state-of-charge region, where calendar ageing is greatest). These bands feed the behaviour layer of Section~\ref{sec:behaviour}.

\subsection{Validation of the Schema}

After conversion the validator checks that the required columns exist and are not entirely empty, that numeric columns can be converted to numbers, that state of charge and state of health lie between 0 and 100 when present, that voltage, capacity and resistance are not negative, and that either a cycle or a timestamp exists. The validator reports each failure with its reason. For the CAN and serial paths, the stop for a missing channel is a separate signal-coverage check, which compares the channels a source provides with the channels the estimators need and refuses the run by name, as when a CAN message carries current and state of charge but no voltage or temperature. The serial frames do not carry a dataset name, so the required-column rule applies to the dataset loaders and not to the live telemetry paths.

\section{DATA FLOW AND PROCESSING SEQUENCE}

Every source follows the same eight steps in the same order. The CAN and serial paths run them through the single function \texttt{score\_telemetry\_frame}, and the dataset studies run the same estimators through the study scripts and the batch pipeline.
\begin{enumerate}
	\item \textbf{Acquire.} Read frames from the bus, lines from the serial port or rows from a file.
	\item \textbf{Convert.} Decode signals or fields and convert to the units and sign convention of the schema.
	\item \textbf{Validate.} Apply the schema checks and the capture-level gates, and refuse the capture if they fail.
	\item \textbf{Segment.} Split the stream into discharges by current direction and by gaps in time.
	\item \textbf{Measure.} Integrate current to obtain charge between voltage thresholds, apply the ohmic correction and estimate SOH and resistance.
	\item \textbf{Predict.} Extrapolate the cell's own fade to the end-of-life threshold and attach an error band.
	\item \textbf{Score behaviour.} Compute the stress features, the risk score, the health index and the Guardian explanation.
	\item \textbf{Publish.} Update the twin and the fleet store, and serve the result with its evidence and its refusals.
\end{enumerate}

\section{DISCHARGE SEGMENTATION}
\label{sec:segment}

The estimator receives voltage $V(t)$, current $I(t)$ and time. The cycler's capacity, resistance and cycle summaries are removed before scoring. Current is segmented into charge, discharge and rest phases using a dead band (0.02~A in the CALCE studies, where rest is logged as exactly zero, and 0.5~A as the default of the live telemetry path) and a minimum run length. The charge delivered on discharge $n$ is the trapezoidal integral of $|I|\,dt$. Every discharge is retained: a partial discharge that crosses the window is a full measurement of the window.

\section{VOLTAGE-WINDOW SOH ESTIMATION}
\label{sec:soh-method}

For a window $[V_\mathrm{low}, V_\mathrm{high}]$, the window charge $\Qw(n)$ is the charge delivered while terminal voltage falls from $V_\mathrm{high}$ to $V_\mathrm{low}$ on discharge $n$. It is read by interpolation and requires that the discharge span at least 95\% of the window. With $\mathcal{R}$ the cell's first five measurable discharges,
\begin{equation}
\SOH(n) = \frac{\Qw(n)}{\operatorname{median}_{m\in\mathcal{R}} \Qw(m)} ,
\label{eq:soh}
\end{equation}
and the reported value is the median of the latest five accepted readings. The primary window, 3.90--3.60~V, was fixed before any validation run. The windows 4.05--3.55 and 3.85--3.55~V are reported as sensitivity analyses and not selected between.

Equation~\eqref{eq:soh} assumes the discharge curve scales approximately uniformly with capacity. Loss of lithium inventory shifts the curve, which the ratio partly absorbs. Loss of active material and resistance growth change its shape and position, which it does not. This is the method's principal approximation, and Section~\ref{sec:res-suff} measures how it drifts with age. Fig.~\ref{fig:sohflow} summarises the estimator.

\begin{figure}[htbp]
    \centering
    \resizebox{!}{0.5\textheight}{%
    \begin{tikzpicture}[box/.style={draw, rounded corners=2pt, align=center, minimum height=9mm, text width=5.4cm, font=\small},
                        gate/.style={draw, rounded corners=2pt, align=center, minimum height=9mm, text width=5.4cm, font=\small, dashed},
                        arr/.style={-{Stealth[length=2.2mm]}, thick}]
        \node[box] (a) at (0,0) {Voltage, current, time};
        \node[box] (b) at (0,-1.7) {Segment discharges; integrate $|I|\,dt$};
        \node[box] (c) at (0,-3.4) {Load-step resistance $R_\mathrm{step}$; shift window by $|I|\tilde R$};
        \node[box] (d) at (0,-5.1) {Window charge $\Qw(n)$ by interpolation};
        \node[gate] (e) at (0,-6.8) {Gates: truncated discharge, plausibility, reference timing};
        \node[box] (f) at (0,-8.5) {$\SOH = \Qw / \operatorname{median}\Qw(\mathcal{R})$};
        \node[box] (g) at (0,-10.2) {Confidence from consistency $u$; error band or refusal};
        \foreach \s/\t in {a/b,b/c,c/d,d/e,e/f,f/g} \draw[arr] (\s.south) -- (\t.north);
    \end{tikzpicture}}
    \caption{Field SOH estimator: from raw telemetry to a banded or refused output.}
    \label{fig:sohflow}
\end{figure}

\section{INTERNAL RESISTANCE ESTIMATION AND OHMIC CORRECTION}
\label{sec:rstep}

Under load, terminal voltage lies below open-circuit voltage by approximately $IR$. As $R$ grows with age, a fixed terminal-voltage window slides along the curve, and the slide reads as additional fade. At each discharge onset, between the last rest sample and the first loaded sample,
\begin{equation}
R_\mathrm{step} = \frac{V_\mathrm{rest} - V_\mathrm{load}}{|I_\mathrm{load} - I_\mathrm{rest}|},
\label{eq:rstep}
\end{equation}
accepted when the two samples are within 60~s and the current change is at least 0.05~A. $R_\mathrm{step}$ is an \emph{apparent} resistance: the ohmic drop plus polarisation accrued by the first loaded sample. On cell CS2\_35 its starting value is about 0.141~$\Omega$ against the cycler's about 0.089~$\Omega$ (reports/metrics/calce\_resistance). Each discharge uses the trailing median $\tilde R(n)$ of the last 11 step estimates, so no correction depends on later data, and its window is shifted to
\begin{equation}
[V_\mathrm{low} - |I_n|\tilde R(n),\; V_\mathrm{high} - |I_n|\tilde R(n)] .
\label{eq:shift}
\end{equation}
Discharges preceding the first resistance estimate are left unmeasured rather than scored uncompensated.

\section{DATA-QUALITY GATES AND REFUSAL CONDITIONS}
\label{sec:gates}

Each gate needs no ground truth and can therefore run in deployment.

\begin{enumerate}
    \item \textbf{Truncated discharge.} A discharge delivering less than 50\% of the reference discharges' total charge is not scored. A cut-short cycle crosses the window with little charge and would read as severe fade.
    \item \textbf{Physical plausibility.} A ratio below 0.50 is withheld, and a cell with more than 20\% of readings below it is refused entirely. A ratio above 1.10 is withheld as not spanning the same part of the curve.
    \item \textbf{Reference timing.} The reference must be complete within 20 equivalent full cycles of charge throughput. Throughput rather than a discharge count is used because storage-protocol records contain thousands of short pulses.
\end{enumerate}

\section{LEARNED VOLTAGE WINDOWS FOR PARTIAL DISCHARGES}
\label{sec:learned}

A car rarely discharges across a fixed window. When the fixed window cannot be read and the rated capacity $C_\mathrm{rated}$ is known, a window is learned from the cell's first ten discharges only. The band they all cover is trimmed by 15\% at each end, and the flattest sub-window of up to 0.30~V within it is selected by its steepness
\begin{equation}
s = \frac{V_\mathrm{high}-V_\mathrm{low}}{\Qw/C_\mathrm{rated}}
\quad [\mathrm{V\ per\ unit\ state\ of\ charge}] .
\label{eq:steep}
\end{equation}
The window is refused if even the flattest choice has $s > 2.0$, because it then lies on the end-of-discharge knee. There, voltage reflects kinetics (resistance and solid-state diffusion) rather than stored charge.

\section{REMAINING USEFUL LIFE ESTIMATION FROM CELL-SPECIFIC FADE}
\label{sec:rul-method}

At cycle $k$, the cell's SOH history up to $k$ is smoothed with an 11-point rolling median, a straight line is fitted to the full history, and its crossing of 0.90 gives the end-of-life estimate. Estimates are refused in three cases: with fewer than 30 points, for flat or rising trends, and when the estimate reaches more than twice the history length ahead. A cell already past the threshold is reported as zero. Linear, square-root and quadratic forms, over full and trailing windows, were compared, and the full-history line was most accurate near end of life. No other cell's data enters, so there is no training cohort.

\section{EVIDENCE SUFFICIENCY CRITERION}
\label{sec:suff-method}

Each output has its own requirements: capacity health needs voltage, current and five reference readings plus one; resistance health needs a clean load step; RUL needs 30 complete discharges; and heat exposure needs a temperature channel. Within capacity health, confidence follows consistency. Because a healthy estimate keeps changing as the cell ages, stability of the estimate is the wrong criterion. The question is whether readings agree once that trend is removed. With $\sigma_{\tilde x}(x_{1..n}) = 1.2533\,\mathrm{MAD}(x)/\sqrt{n}$ the standard error of a median,
\begin{align}
u_\mathrm{ref} &= \frac{\sigma_{\tilde x}\!\left(\Qw(\mathcal{R})\right)}{\operatorname{median}\Qw(\mathcal{R})}, \\
u_\mathrm{rec} &= \frac{1.2533\,\mathrm{MAD}(\varepsilon)}{\sqrt{5}}, \\
u &= \sqrt{u_\mathrm{ref}^2 + u_\mathrm{rec}^2},
\label{eq:u}
\end{align}
where $\varepsilon$ are the residuals of the last ten readings about their own straight line. The tolerance $u \le 0.01$ (one SOH point) was fixed before the study of Section~\ref{sec:res-suff}, chosen to lie below the validated 1.7\% median error. Confidence is then assigned as \emph{high} (consistent, fixed window and a beginning-of-life reference), \emph{medium} (consistent, but a learned window or a start-of-log reference), or \emph{low} (inconsistent, reported with the measured 22-point band).

\section{HEURISTIC BEHAVIOUR SCORING LAYER}
\label{sec:behaviour}

Alongside the measured estimators, BEACON carries a rule-based behaviour layer inherited from the project's first phase. It converts usage features into a risk score, an aging budget, a health index, a state label, a Guardian explanation and a digital-twin state. The features are average stress, average temperature, deep-discharge duration, fast-charge duration, aggressive-discharge count and average SOC.

\begin{table}[htbp]
	\caption{Risk Categories of the Rule-Based Layer}
	\centering
	\small
	\renewcommand{\arraystretch}{1.15}
	\setlength{\tabcolsep}{4pt}
	\begin{tabular}{|C{0.25\linewidth}|C{0.25\linewidth}|}
		\hline
		\textbf{Risk category} & \textbf{Risk score} \\
		\hline
		LOW & 0--39 \\
		\hline
		MEDIUM & 40--59 \\
		\hline
		HIGH & 60--79 \\
		\hline
		CRITICAL & 80--100 \\
		\hline
	\end{tabular}
	\label{tab:risk}
\end{table}

\begin{table}[htbp]
	\caption{Health Index, Battery State and Digital-Twin State}
	\centering
	\small
	\renewcommand{\arraystretch}{1.15}
	\setlength{\tabcolsep}{4pt}
	\begin{tabular}{|C{0.22\linewidth}|C{0.22\linewidth}|C{0.28\linewidth}|}
		\hline
		\textbf{Health index} & \textbf{Battery state} & \textbf{Twin state} \\
		\hline
		0--29 & HEALTHY & NORMAL \\
		\hline
		30--59 & WARNING & MODERATE\_RISK \\
		\hline
		60--79 & DEGRADED & HIGH\_RISK \\
		\hline
		80--100 & CRITICAL & FAILURE\_IMMINENT \\
		\hline
	\end{tabular}
	\label{tab:rules}
\end{table}

The aging budget starts at 100 and is reduced by penalties for stress, temperature, deep discharge, fast charging, aggressive discharge and SOC abuse, and the health index is $\mathrm{HI} = 100 - \text{budget}$, so a higher index means a worse battery (Tables~\ref{tab:risk} and~\ref{tab:rules}). The digital twin is a one-to-one relabelling of the battery states, with a queryable API. It adds structure but no new predictive model.

This layer is \emph{not} validated against measured degradation. The Spearman correlation between the risk score and measured fade on a separate set of 33 NASA cells from the earlier health-index ranking was $\rho = -0.27$ ($p = 0.12$), so it is labelled a heuristic wherever it appears. Where measured SOH exists and the log starts at beginning of life, the measured SOH sets the battery state and records that it did, and the heuristic index is used only otherwise.

\section{DIGITAL-TWIN STATE MODEL}

The digital twin records, for each battery, a state derived from the health index of the latest run and the history of changes in that state. The state is a relabelling of the health index (Table~\ref{tab:rules}), not a separately tuned threshold set, so any change to the health-index formula flows through to the twin without a second edit. A transition is recorded whenever the state of a battery differs from its state in the previous run, and a battery's first observation counts as a transition from no state. Re-running an unchanged, identically seeded fleet therefore produces no new transitions, which is checked by a test. The failure likelihood the twin reports is a monotonic transform of the health index onto the interval from 0 to 1. It is named a likelihood and not a probability because it is not a calibrated statistical quantity.

\section{DESIGN DECISION RECORDS}

Fourteen architecture decision records document the choices made during the project, with the context, the decision and its status; the table gives a short title for each. Table~\ref{tab:adr} lists them. Several record a claim that was withdrawn after a later measurement, and that history is kept in the repository and not rewritten.

\begin{table}[htbp]
	\caption{Architecture Decision Records}
	\centering
	\footnotesize
	\renewcommand{\arraystretch}{1.1}
	\setlength{\tabcolsep}{3pt}
	\begin{tabular}{|C{0.06\linewidth}|L{0.8\linewidth}|}
		\hline
		\textbf{No.} & \textbf{Decision} \\
		\hline
		1 & Cross-dataset validation is not possible; leave-one-cohort-out replaces it \\
		\hline
		2 & Health index version 1 stays the default; version 2 is exposed conditionally \\
		\hline
		3 & Exact Shapley attribution for the rule scores; SHAP only where it is the right tool \\
		\hline
		4 & The dashboard renders only computed values \\
		\hline
		5 & The gate is the product, not the model \\
		\hline
		6 & Transfer targets are screened on commensurability, not availability \\
		\hline
		7 & The target was the binding constraint, not the model \\
		\hline
		8 & Model capacity buys within-protocol skill and spends it on transfer \\
		\hline
		9 & The protocol-shift collapse replicates on CALCE, but its shape is dataset-specific \\
		\hline
		10 & The leave-one-cohort-out estimate is noisier than the effect it is used to measure \\
		\hline
		11 & Unsupervised detection finds structure, not harm \\
		\hline
		12 & An absolute SOH target on CALCE recovers four cells and refutes the claim that learned models add enormously \\
		\hline
		13 & The estimator-variance result survives the target correction as a ratio \\
		\hline
		14 & Rated capacity is declared, not defaulted \\
		\hline
	\end{tabular}
	\label{tab:adr}
\end{table}

\section{ALGORITHM SUMMARY}

The estimators of the previous sections combine into one procedure per cell, summarised below in the order in which the steps run. Every step that can fail ends in a refusal with a stated cause and not in a default.

\begin{footnotesize}
\begin{verbatim}
for each cell in the capture:
    discharges <- segment(stream)           # by current, time gaps
    if capture-level gate fails:  refuse(cause)
    for each discharge n:
        R(n)  <- trailing median of last 11 load-step estimates
        if no R yet:  mark n "unmeasured"; continue
        window <- [V_low, V_high] shifted by |I_n| * R(n)
        if window not spanned:  mark n "unmeasured"
        if truncated / implausible / reference too short:  refuse n
        Q(n)  <- charge delivered inside window (coulomb count)
        SOH(n) <- Q(n) / Q_ref
    if enough complete discharges (30):
        RUL <- cycles until fade of complete-discharge capacity
               crosses the EOL threshold (90 % default)
    else:  refuse RUL
    if evidence inconsistent:  confidence <- low
    return SOH, R, RUL, confidence, band, refusals
\end{verbatim}
\end{footnotesize}

Remaining life is taken from the capacity of complete discharges only: extrapolating the windowed health of partial discharges was tested and is not accurate enough to report (Chapter~\ref{ch:results}). The reference $Q_{\mathrm{ref}}$ is the declared rated capacity when the source gives one, or the cell's own early capacity otherwise, and the source of the reference is returned with the result. The error band attached to each output is the measured error of that output on the validation data and is not a quantity computed from the cell being scored.

\section{VALIDATION PROTOCOL}
\label{sec:protocol}

\paragraph{Ground Truth.} Ground truth is the cycler's measured discharge capacity, which never enters an estimate. SOH truth is capacity over the median of the cell's first five cycles. For partial-cycling cells, whose checks are about 100 cycles apart, the partial-discharge study scores each cell against its first check, and the remaining-life study uses the capacity over the checks within the first ten cycles. Each is a reference a BMS can hold. A reference drawn from the cell's whole record uses the future, and Section~\ref{sec:res-rul} quantifies the resulting inflation.

\paragraph{Pre-Declared Configurations.} Primary windows, comparison arms, tolerances and scoring were fixed before each study ran. Choices made after seeing a result are labelled exploratory.

\paragraph{Benchmark.} The benchmark uses leave-one-cell-out (LOBO, 31 folds) and leave-one-cohort-out (LOCO, 9 folds) validation. $R^2$ is computed against the training-fold mean, with 95\% bootstrap intervals resampling folds (2,000 resamples). Hyperparameters were fixed and not tuned per fold, except ElasticNet's penalty, chosen by inner cross-validation within the training fold.

\paragraph{Aggregation.} Errors are summarised per cell, then across cells, so long-lived cells do not dominate.

\section{PARAMETERS OF THE ESTIMATORS}

The estimators have few parameters, and each was fixed on the development data before the external datasets were scored. Table~\ref{tab:params} lists those quoted in this chapter. None was changed when the estimators were applied to the Oxford and NASA data.

\begin{table}[htbp]
	\caption{Parameters Fixed Before External Validation}
	\centering
	\small
	\renewcommand{\arraystretch}{1.15}
	\setlength{\tabcolsep}{4pt}
	\begin{tabular}{|L{0.38\linewidth}|C{0.14\linewidth}|L{0.28\linewidth}|}
		\hline
		\textbf{Parameter} & \textbf{Value} & \textbf{Used for} \\
		\hline
		Maximum gap between rest and load samples & 60~s & Accepting a load-step estimate \\
		\hline
		Minimum current change at the step & 0.05~A & Accepting a load-step estimate \\
		\hline
		Length of trailing resistance median & 11 steps & Ohmic correction \\
		\hline
		Rest threshold of mode inference & $\pm$0.02~A & Schema-level operating mode \\
		\hline
		Serial sampling period (rig) & 1~s & Rig firmware \\
		\hline
		Default end-of-life threshold & 90\% & RUL target (80\% is the convention but is seldom reached on CALCE) \\
		\hline
	\end{tabular}
	\label{tab:params}
\end{table}

Fixing the parameters in advance has a cost. They are not the best possible values for any one dataset, and the errors reported on the external data are therefore a fair statement of what the estimators do on data they were not tuned on, and not the best that could be obtained by tuning.

\section{EVALUATION METRICS}
\label{sec:metrics}

\begin{itemize}
    \item \textbf{SOH error:} per-cell mean absolute error in SOH (0.01 equals one percentage point), the median over cells, a 95\% interval that resamples cells, and the worst cell.
    \item \textbf{Coverage:} cells measured out of cells scored. Refusals count against coverage and are reported with their cause.
    \item \textbf{Resistance:} Spearman correlation against the cycler's resistance column over each cell's life, direction of growth, and growth ratio.
    \item \textbf{RUL:} the share of estimates within $\pm$20 cycles, the median error, and the median bias, by true distance to end of life and by predicted distance.
    \item \textbf{Sufficiency:} median and 90th-percentile error by consistency class, and the correlation between $u$ and absolute error.
    \item \textbf{Benchmark:} LOBO and LOCO $R^2$ against the training mean with bootstrap intervals, and LOCO mean absolute error.
\end{itemize}

\section{EXPERIMENTAL CONTROLS AND REPRODUCIBILITY}
\label{sec:repro}

Every reported number is regenerated by a named script from committed artifacts (Appendices~A and~B), and a test (\texttt{test\_reported\_numbers}) ties each figure quoted in the documentation to the artifact that produced it, so a re-run experiment with a stale paragraph fails the build. A control shuffles each cell's capacity in time to confirm that the benchmark harness does not manufacture skill. A continuous-integration job reproduces the headline benchmark end to end.
